%% Power Shield: AYLA Pipeline — Condensed 6-Page Journal Submission
%% Compile with: pdflatex -> bibtex -> pdflatex -> pdflatex

\documentclass[twocolumn, 10pt, letterpaper]{article}

\usepackage[top=0.75in, bottom=0.75in, left=0.65in, right=0.65in]{geometry}
\usepackage{amsmath, amssymb}
\usepackage{graphicx}
\usepackage{booktabs}
\usepackage{hyperref}
\usepackage{cite}
\usepackage{array}
\usepackage{caption}
\usepackage{float}
\usepackage{microtype}
\setlength{\columnsep}{18pt}

\hypersetup{
    colorlinks=true,
    linkcolor=blue,
    citecolor=blue,
    urlcolor=blue
}

\title{\textbf{Power Shield:}\\
An Automated Yield-Loss Alerting (AYLA) Pipeline\\
Using Explainable AI (XAI) and Machine Learning (ML)\\
for Power Grid Attack Detection and LLM-Guided Remediation}

\author{
    Gabriel Adams\\
    \textit{Department of Computer Science}\\
    \textit{Tennessee Technological University}
}

\date{May 2026}

\begin{document}

\twocolumn[{%
  \maketitle
  \begin{abstract}
  Power grids face coordinated cyber-physical attacks---including False Data Injection
  (FDIA), Relay Saturation, and Remote Terminal Communication Injection---designed to
  evade conventional anomaly detectors by mimicking plausible physical behavior.
  This paper presents \textbf{Power Shield}, implementing an
  \textbf{Automated Yield-Loss Alerting (AYLA)} pipeline: a four-stage explainable AI
  system that detects grid attacks, attributes each alert to specific sensor drivers via
  SHAP, retrieves the matching remediation protocol via RAG, and delivers step-by-step
  operator guidance via a local LLM.
  The SHAP-guided stacked ensemble achieves \textbf{91.88\%} accuracy and
  \textbf{94.6\%} attack recall on the full combined MSU/ORNL SCADA dataset
  (78,377 samples), surpassing the established combined-dataset baseline.
  Stages 2--4 (SHAP, RAG, and LLM) activate only on classifier-flagged events,
  not on every sensor reading, so the expensive LLM inference is incurred only
  when explanation and remediation are genuinely needed.
  A no-context ablation confirms LLM grounding: withholding retrieved protocol
  documents drops attribution accuracy from \textbf{90.5\%} to \textbf{16.5\%}
  ($-74.0$~pp, non-overlapping 95\% CIs), proving that pipeline accuracy reflects
  genuine document grounding rather than parametric knowledge.
  \end{abstract}
  \vspace{1em}
}]

% ─────────────────────────────────────────────────────────────────────────────
\section{Introduction}
\label{sec:intro}
% ─────────────────────────────────────────────────────────────────────────────

Power grids are among the most critical infrastructure systems in modern society.
State-sponsored attacks---including the Ukrainian blackouts of 2015--2016
\cite{liang2017ukraine}---demonstrated that adversaries can weaponize SCADA control
software with devastating physical consequences.
ICS attacks increased by over 140\% between 2020 and 2024.

Traditional SCADA intrusion detection relies on rule-based signature matching that
fails against False Data Injection Attacks (FDIA), which craft a sensor manipulation
vector that passes the bad-data detection residual test \cite{liu2011fdia}.
Machine learning classifiers have emerged as a complement to signature-based systems,
learning multi-variate anomaly patterns directly from sensor data.
However, detection alone is insufficient: an operator who receives an alert without
knowing \textit{which} sensors triggered it, \textit{which} attack scenario it matches,
and \textit{what} corrective steps to take is at best slowed and at worst misled.
Current commercial ICS/OT platforms (Claroty, Dragos, Nozomi Networks) surface alerts
but do not generate automated, grounded, step-by-step remediation guidance.

A critical methodological distinction separates prior literature on the
MSU/ORNL power grid dataset: prior work achieving high detection accuracy
\cite{panthi2022ensemble,naeem2025deep,graphkan2025} evaluates classifiers on
individual attack sub-scenarios \textit{in isolation}, eliminating the inter-class
ambiguity that characterises combined-dataset deployment.
The only prior work applying explainability to the \textit{combined} MSU/ORNL dataset
\cite{shapgrid2026} demonstrates SHAP feature attribution but does not extend
to automated protocol retrieval or operator remediation.

\textbf{Power Shield} implements AYLA to close this gap end-to-end.
Its three principal contributions are:

\begin{enumerate}
    \item[\textbf{C1.}] \textbf{SHAP-guided RAG query construction.}
          Per-event SHAP attributions identify the top contributing sensors
          and enrich their names with semantic type descriptors
          (e.g., \texttt{R1-PA:ZH} $\rightarrow$ \texttt{R1-PA:ZH (impedance harmonic)})
          before forming the RAG retrieval query, directly linking the classifier's
          explanation layer to document retrieval.
          Prior work uses SHAP for post-hoc feature explanation
          \cite{shap_ids,shapgrid2026};
          AYLA uses per-event attributions as \textit{semantic retrieval keys}---to
          the author's knowledge, a novel application.

    \item[\textbf{C2.}] \textbf{Ablation-confirmed LLM grounding (XAI transparency verification).}
          A no-context ablation withholds retrieved protocol documents while keeping
          all other prompt inputs (SHAP drivers, sensor readings, classifier verdict)
          identical.
          Attribution accuracy drops from 90.5\% to 16.5\%---a 74.0~pp decline
          with non-overlapping 95\% CIs---directly confirming that full-pipeline
          accuracy reflects genuine document grounding, not parametric knowledge.
          This ablation is an XAI-style transparency verification: it proves
          empirically that the pipeline's recommendations trace to specific
          retrieved protocol documents rather than opaque parametric memory,
          establishing end-to-end decision traceability from sensor reading
          through SHAP attribution, retrieval, and LLM reasoning.

    \item[\textbf{C3.}] \textbf{Combined-dataset evaluation surpassing the baseline.}
          AYLA operates on the full heterogeneous MSU/ORNL dataset and surpasses the
          Hink et al.\ \cite{hink2014ml} combined-dataset baseline (90.56\%)---the only
          prior result directly comparable to a combined-dataset formulation---while
          adding SHAP attribution, RAG protocol retrieval, and LLM-guided remediation.
\end{enumerate}

% ─────────────────────────────────────────────────────────────────────────────
\section{System Architecture}
\label{sec:system}
% ─────────────────────────────────────────────────────────────────────────────

The AYLA pipeline processes each grid event through four sequential stages.
Stages 1--2 operate on raw sensor data; Stages 3--4 operate on the symbolic output
of Stage 2.
Critically, Stages 2--4 (SHAP attribution, RAG retrieval, and LLM remediation)
activate only when Stage 1 exceeds $\tau$: the classifier runs at machine speed on
every incoming sensor snapshot, while LLM inference is reserved for the small
fraction of events that require explanation and remediation.
This event-driven design makes the pipeline viable for real-time SCADA deployment
without high-throughput inference infrastructure.

\medskip\noindent
\textbf{Stage 1 --- Probabilistic Classification.}
A LightGBM/XGBoost stacked ensemble trained on 29 SHAP-selected features
(from an initial pool of 128) ingests a sensor snapshot and outputs
$P(\text{Attack})$ via a LogisticRegression meta-learner.
If $P(\text{Attack}) > \tau = 0.560$, the event is flagged and passed to Stage~2.

\medskip\noindent
\textbf{Stage 2 --- SHAP Attribution.}
TreeSHAP \cite{lundberg2020treeshap} computes exact SHAP values for all 128 features
via the LightGBM base model.
The top-$k$ positive contributors (sensors pushing toward Attack) are selected
as \textit{alert drivers} and expanded to semantic descriptors using a relay-and-type
regex registry.
SCADA, relay, and Snort IDS log channels are reported independently from physical
sensor drivers.

\medskip\noindent
\textbf{Stage 3 --- RAG Protocol Retrieval.}
A two-phase retrieval strategy queries a ChromaDB vector store indexed from
thirteen synthetic grid-protocol manuals using \texttt{all-MiniLM-L6-v2}
sentence embeddings.
\textit{Phase~1} (bi-encoder, multi-pass): a relay-filtered primary query ($k{=}4$)
surfaces relay-specific FDIA protocols; supplementary type-pattern queries ($k{=}2$)
surface generic protocols.
\textit{Phase~2} (cross-encoder reranking, adopted from IHGR-RAG
\cite{ihgrrag2025}): all candidates are rescored jointly against the primary query
by a \texttt{ms-marco-MiniLM-L-6-v2} cross-encoder, producing scores comparable
across all retrieval passes regardless of which pass found each candidate.

\medskip\noindent
\textbf{Stage 4 --- LLM Remediation.}
The ranked protocol context, SHAP summary, sensor readings, and classifier verdict
are assembled into a structured prompt.
A Chain-of-Thought (CoT) \cite{wei2022cot} step instructs the LLM to explicitly
verify the RANK~\#1 protocol's diagnostic signature against the observed sensors
before committing to a diagnosis.
The model outputs: (1)~Protocol Match Reasoning, (2)~Diagnosis with protocol ID,
and (3)~numbered remediation steps drawn directly from the matched manual.

\begin{table*}[t]
\centering
\caption{AYLA system component configuration.}
\label{tab:config}
\small
\begin{tabular}{ll}
\toprule
\textbf{Component} & \textbf{Configuration} \\
\midrule
Classifier        & LightGBM \cite{ke2017lightgbm} + XGBoost \cite{chen2016xgboost}
                    ensemble (29/128 SHAP-selected features) \\
Meta-learner      & LogisticRegression ($\beta_1{=}6.74$, $\beta_2{=}0.68$) \\
Alert threshold   & $\tau = 0.560$ (threshold-sweep optimised; false alert $< 15\%$) \\
SHAP              & TreeSHAP \cite{lundberg2020treeshap} (exact; LightGBM base model;
                    zero-padded to 128-d) \\
Embedding model   & \texttt{all-MiniLM-L6-v2} (384-dim); ChromaDB (L2 distance) \\
RAG Phase 1       & Bi-encoder; $k{=}4$ relay-filtered + $k{=}2$ type-pattern passes \\
RAG Phase 2       & Cross-encoder reranking (\texttt{ms-marco-MiniLM-L-6-v2}) \\
LLM               & Mistral-Nemo 12B (Ollama, local; temperature 0.0; max 2048 tokens) \\
\bottomrule
\end{tabular}
\end{table*}

% ─────────────────────────────────────────────────────────────────────────────
\section{Mathematical Formulation}
\label{sec:math}
% ─────────────────────────────────────────────────────────────────────────────

\textbf{SHAP-guided feature selection.}
An initial LightGBM model trained on all $d{=}128$ features computes global
mean absolute SHAP importance on a 3000-sample hold-out.
Features whose mean absolute SHAP value meets or exceeds the mean across all
128 features are retained, reducing $d{=}128$ to $d'{=}29$ features.

\textbf{Ensemble stacking.}
A LogisticRegression meta-learner combines LightGBM and XGBoost
attack probabilities on a held-out calibration set:
\begin{equation}
    P(\text{Attack} \mid \mathbf{x}) =
    \sigma\!\left(\beta_0
    + \beta_1 P_{\text{LGB}}(\mathbf{x}_S)
    + \beta_2 P_{\text{XGB}}(\mathbf{x}_S)\right)
    \label{eq:stack}
\end{equation}
where $\mathbf{x}_S$ is the 29-feature SHAP-selected subset,
$\beta_1 = 6.74$, $\beta_2 = 0.68$, and $\tau = 0.560$ is selected by
sweeping $[0.30, 0.90]$ subject to false alert rate $< 15\%$.

\textbf{Cross-encoder reranking.}
For primary query $q$ and candidate document $d_j$, the cross-encoder
$f_{\text{CE}}$ scores each pair jointly:
\begin{equation}
    \hat{s}_j = f_{\text{CE}}(q, d_j) + f_{\text{CE}}(q_{\text{freq}}, d_j)
    \label{eq:ce}
\end{equation}
where $q_{\text{freq}}$ is added only when frequency sensors
(\texttt{:F}, \texttt{:DF}) dominate the alert drivers, boosting FREQ-01
retrieval without penalising relay-specific FDIA cases.
Unlike bi-encoder L2 scores, cross-encoder scores are directly comparable
across all retrieval passes.

% ─────────────────────────────────────────────────────────────────────────────
\section{Dataset and Experimental Setup}
\label{sec:setup}
% ─────────────────────────────────────────────────────────────────────────────

\subsection{Dataset}

Experiments use the MSU/ORNL ICS Security Dataset \cite{morris2014ics},
a widely used power system intrusion detection benchmark generated on a
hardware-in-the-loop testbed with a 4-relay transmission protection system
communicating via DNP3 and Modbus over a simulated SCADA network.
The binary variant contains \textbf{78,377 samples} with \textbf{128 features}:
voltage and current magnitudes, impedance and phase harmonics, relay-level
frequency measurements, and SCADA/relay/Snort IDS log channels, spanning
four protection relays (R1--R4).
The class distribution is 55,663 Attack (71.0\%) and 22,714 Natural (29.0\%)
(2.45:1 imbalance, corrected via \texttt{is\_unbalance=True} and
\texttt{scale\_pos\_weight}).
Data is split 64/16/20 into training, calibration, and held-out test sets
($n{=}15{,}676$ test samples).

\subsection{Protocol Knowledge Base}

Thirteen protocol documents were authored by the researcher, grounded in the
MSU/ORNL dataset's documented attack signatures \cite{morris2014ics} and
standard power system protection principles.
Each document specifies the distinguishing sensor pattern for one attack or fault
class and a set of remediation steps consistent with conventional substation
procedures; they serve as a proxy for operational runbooks that would be
indexed in a production deployment.
The documents are indexed into ChromaDB using 700-character chunks with
100-character overlap and span relay-specific FDIA scenarios (R1-FDIA through
R4-PM), generic anomaly classes (RSCA-01, RTCI-01, FREQ-01), fault and
degradation scenarios (CURR-01, PHASE-01, FAULT-01, DEGR-01, DEGR-02), and a
normal-operation baseline (EXPECTED-09).
The following excerpt from the R1-FDIA entry illustrates the protocol format:

\begin{quote}\small
\textbf{R1-FDIA:} If voltage magnitude sensors on Relay~1 (\texttt{R1-PM1:V},
\texttt{R1-PM2:V}, \texttt{R1-PM3:V}) show fluctuations exceeding 10\% while the
corresponding impedance sensor (\texttt{R1-PA:Z}) remains flat and no current spike
is observed on \texttt{R1-PM4:I}--\texttt{R1-PM6:I}, the likely cause is a False
Data Injection Attack targeting Relay~1 phasor measurements.
\textit{Remediation:} Lock Relay~1 data stream from SCADA write access; manually
verify phasor data at the R1 substation; if discrepancy with R2 persists beyond
30~s, escalate to control center and flag for cybersecurity review.
\end{quote}

\subsection{Evaluation Protocol}

\textbf{Classifier evaluation} uses the full 15,676-sample held-out test set.

\textbf{RAG retrieval accuracy} is assessed on 13 hand-crafted test cases
(one per protocol); a case is a hit if the expected protocol identifier appears
anywhere in the retrieved context.

\textbf{LLM attribution accuracy} is evaluated on a stratified random sample of
200 triggered events.
Ground truth is assigned by matching top SHAP sensors against curated protocol
sensor lists via a weighted overlap score; a response is correct if the expected
protocol ID appears anywhere in the LLM output.

\textbf{No-context ablation} re-evaluates the same 200 samples with the retrieved
protocol context replaced by a placeholder, holding all other prompt inputs
constant.

% ─────────────────────────────────────────────────────────────────────────────
\section{Results}
\label{sec:results}
% ─────────────────────────────────────────────────────────────────────────────

\subsection{Classifier Performance}

\begin{table}[H]
\centering
\caption{AYLA ensemble classifier on the held-out 20\% test set
         ($n{=}15{,}676$). CIs are Wilson 95\%.}
\label{tab:classifier}
\small
\begin{tabular}{lrr}
\toprule
\textbf{Metric} & \textbf{Value} & \textbf{95\% CI} \\
\midrule
Overall Accuracy   & 91.88\% & [91.44\%, 92.30\%] \\
Attack Recall      & 94.55\% & [94.11\%, 94.95\%] \\
Natural Recall     & 85.34\% & [84.28\%, 86.34\%] \\
False Alert Rate   & 14.66\% & [13.66\%, 15.72\%] \\
AUC-ROC / MCC      & 0.9665  & 0.8020 \\
\bottomrule
\end{tabular}
\end{table}

The ensemble achieves AUC-ROC of 0.9665 and MCC of 0.8020, indicating strong
discriminative power on the full combined dataset despite the 2.45:1 class imbalance.
AYLA's 91.88\% accuracy surpasses the Hink et al.\ \cite{hink2014ml}
combined-dataset baseline of 90.56\%---the only directly comparable prior result.
Prior work reporting 95--99\% accuracy \cite{panthi2022ensemble,naeem2025deep,graphkan2025}
evaluates on isolated per-scenario sub-datasets, eliminating inter-class ambiguity
and making those results not directly comparable.

\subsection{RAG Retrieval Accuracy}

The two-phase retrieval achieves perfect recall on all 13 curated protocol test
cases: 13/13 (100\%) of expected protocol identifiers appear in the retrieved context.
Cross-encoder reranking is critical for several protocols (R4-PM, RTCI-01, FREQ-01)
where the correct document is retrieved but does not rank first from the bi-encoder alone.

\subsection{LLM Protocol Attribution and Ablation}

\begin{table}[H]
\centering
\caption{LLM protocol attribution accuracy and no-context ablation
         (same 200 samples, identical prompt structure). 95\% Wilson CIs shown.}
\label{tab:llm_ablation}
\footnotesize
\setlength{\tabcolsep}{4pt}
\begin{tabular}{lrr}
\toprule
\textbf{Condition} & \textbf{Correct} & \textbf{95\% CI} \\
\midrule
Full pipeline (RAG context) & 181/200 (90.5\%) & [85.6\%, 93.8\%] \\
No-context ablation         &  33/200 (16.5\%) & [12.0\%, 22.3\%] \\
\midrule
$\chi^2$ ($p_0{=}1/13$) & \multicolumn{2}{r}{1931.42, $p < 10^{-10}$} \\
Accuracy drop & \multicolumn{2}{r}{$-74.0$~pp} \\
\bottomrule
\end{tabular}
\end{table}

The full AYLA pipeline achieves 90.5\% correct protocol citation (181/200) with a
100\% any-citation rate.
A $\chi^2$ goodness-of-fit test ($\chi^2 = 1931.42$, $p < 10^{-10}$) confirms
performance far exceeds the random-selection baseline ($p_0 = 1/13 \approx 7.7\%$).

Removing retrieved context drops attribution to 16.5\%---a \textbf{74.0~pp decline}
with non-overlapping confidence intervals
([85.6\%, 93.8\%] vs.\ [12.0\%, 22.3\%]).
Without context, the LLM collapses to citing a single default protocol (R4-PM)
for nearly all events regardless of relay identity or sensor type, consistent with
parametric guessing rather than document-grounded reasoning.
This ablation directly confirms that the 90.5\% pipeline accuracy is grounded in
the retrieved protocol manuals.

% ─────────────────────────────────────────────────────────────────────────────
\section{Discussion}
\label{sec:discussion}
% ─────────────────────────────────────────────────────────────────────────────

\textbf{Remaining failure mode: RTCI-01.}
RTCI-01 (Remote Terminal Communication Injection) is the dominant unsolved failure
class (0/7 = 0\%).
RTCI-01 involves current collapse correlated with Snort IDS log spikes; however,
SHAP attribution frequently elevates phase-angle or current-harmonic sensors above
the log channels when both co-occur, depriving the RAG query of the vocabulary
needed to distinguish RTCI-01 from relay-specific FDIA scenarios.
A log-channel promotion fix was attempted---elevating Snort log drivers in
the RAG query regardless of SHAP rank when active---but this introduced new
failures in other protocol classes: Snort log activity co-occurs across multiple
attack types, so unconditionally promoting log channels over-retrieves RTCI-01
for events that are actually FDIA or FREQ-01.
RTCI-01 remains an open failure class; resolving it likely requires a dedicated
log-channel detection stage or protocol-aware retrieval filtering rather than a
SHAP-rank override.
FREQ-01 improved from 0\% in the original single-LightGBM system to 75.0\% with
cross-encoder reranking, confirming that score incomparability across retrieval
passes was the barrier, not missing protocol coverage.

\textbf{Protocol knowledge base realism.}
The 13 protocol documents were authored by the researcher, grounded in the
MSU/ORNL dataset's documented attack signatures and standard power system
protection principles, but not independently validated by working grid operators.
In production, the vector store would be indexed from actual operator runbooks,
NERC CIP procedures, or vendor relay documentation.
The pipeline is protocol-agnostic: the indexed document set is a deployment
parameter, not an architectural constraint, and results here represent a
proof-of-concept on synthetic protocols.

\textbf{Hardware and deployment.}
All results were produced on a laptop NVIDIA RTX~3060 (6~GB VRAM), averaging
\textbf{3.5 minutes per event}.
Because the LLM stage is event-driven---triggering only when a sensor snapshot
exceeds $\tau$, not on every relay log---the 3.5-minute cost is incurred only on
flagged anomalies.
For production: an RTX~4090 reduces latency to $\sim$45--90~s per event;
an A100 to $\sim$20--40~s.
For single-substation deployments with rare attack events, laptop-class hardware
is viable; multi-substation control centers should use a dedicated GPU server.
Local LLM deployment (Ollama) is required in environments where SCADA telemetry
cannot be transmitted to external API endpoints.

% ─────────────────────────────────────────────────────────────────────────────
\section{Conclusion}
\label{sec:conclusion}
% ─────────────────────────────────────────────────────────────────────────────

Power Shield implements AYLA, an end-to-end pipeline for power grid attack
detection and LLM-guided remediation.
Three results establish its feasibility.
\textbf{First}, the SHAP-guided LightGBM/XGBoost ensemble achieves 91.88\%
accuracy, AUC-ROC 0.9665, and 94.6\% attack recall on the full combined
MSU/ORNL dataset, surpassing the Hink et al.\ \cite{hink2014ml} combined-dataset
baseline---the only directly comparable prior result.
\textbf{Second}, a no-context ablation confirms that the 90.5\% LLM protocol
attribution (181/200) is grounded in retrieved documents: removing context drops
accuracy to 16.5\%, a 74.0~pp decline with non-overlapping 95\% CIs.
\textbf{Third}, unlike current commercial ICS security platforms (Claroty, Dragos,
Nozomi Networks), which surface alerts but provide no automated remediation,
AYLA delivers step-by-step, protocol-grounded operator guidance for each
flagged event---a capability absent from all known deployed ICS security systems.
The full implementation, pre-trained model, and evaluation data are publicly
available at \url{https://github.com/gradams42/power-shield}.

\clearpage
\bibliographystyle{unsrt}
\bibliography{references}

\end{document}
